% !TeX root = ../papers/nc-p3.tex
\section{有限な第一窓の剛性と三次障害}
\label{sec:finite-window-rigidity}
前節の二十点例で、正規化した第三空間の候補を退化解まで含めて決定する。
これにより第一窓を動かしたときの非延長開集合も得られる。以下では第三階数を指定しない。

\begin{theorem}[二十点例の全平方根]\label{thm:finite-all-square-roots}
前節の被約例の第一座標空間を $U\subset D=k^{20}$、$Q=U^2$とする。
$\Gr(4,D)$で$L^2\subset Q$を課した閉スキームは、ちょうど$2^{19}$個の被約点からなる。
その全点は$L=\epsilon U$で、$\epsilon$は各座標で独立に取った符号、全体符号は同一視する。
従って、この第一窓から四辺の点対応が有限な標準AS対象への延長は添字保持同型を除き一つである。
\end{theorem}
\begin{proof}
二十点を行とする第一座標行列を$X$とする。
四座標関数を$x_1,\ldots,x_4$とし、$x_0=0$と置く。
十個の一次式$a_i$として$x_r-x_s$ ($0\le r<s\le4$)を取り、
その零点の行集合を$P_i$とする。以下の小さい行列計算は
\path{computations/verify_finite_window_rigidity.py}で有理数上再現できる。
\begin{enumerate}
\item 各$P_i$は七行で、$Q$の制限は階数6である。その評価七行の唯一の依存の係数は全て非零である。
$P_i\cup P_j$への制限は$i\ne j$で階数9である。他の九つの一次式は常に$U$を張る。
\item $P_i\cup P_j\cup P_k$の補で$X$が階数4なら三つ組を採用する。
採用されるのは120組中80組である。任意の一対から始め、採用した三つ組の一辺があれば残り二辺も加えると、
どの始点からも全45対に達する。
\item 最初の平面と各他の平面の和集合の補には、どの四行も独立な五行がある。
十平面は共通点を持たない。
\end{enumerate}
$L$の平方が$Q$に入るとする。第一の依存の係数を対角に置けば七次元非退化対称形式を得る。
$L$の七行への制限は全等方なので次元は3以下である。
従って各$i$で$P_i$上零になる$0\ne h_i\in L$を選べる。
$P_i$上零の$Q$の元はちょうど$a_iU$である。これは次元が4で、$U$の二次単項式は独立だからである。
また$P_i\cup P_j$上零の$Q$の元は$ka_i a_j$である。

$h_i h_j\ne0$ならこれは$a_i a_j$の非零スカラー倍である。そのような対が存在する。
全対で零なら十個の$h_i$の非零座標の台が互いに交わらず、四次元空間内で十個が独立になってしまう。
一対$i,j$で積が非零のとき、採用三つ組$i,j,k$を取る。
もし$h_i h_k=0$なら$h_k$は$P_i\cup P_j$の補で零である。
$h_k^2=a_k b$ ($b\in U$)なので、三平面の補で$b$が零になり、階数4から$b=0$となる。
被約性と$h_k\ne0$に反する。同じ議論を$j$にも用いて第二の計算を適用すれば、全対で積が非零となる。

従って各$i$で$h_i L=a_i U$である。右辺の基底を他の$a_j$との積から取れるためである。
$P_i$の補では右辺に基点がなく、$h_i$はそこで消えない。
この等号の与える線型同型から、$L$と$U$のその補上の評価点列は一つの射影同型で対応する。
最初の補との共通部分に射影枠があるので、これらの射影同型は全て等しい。
補が全二十行を覆うため、各行の比例定数をまとめた単元$d$で$L=dU$となる。
$d^2Q\subset Q$なら$d^2$はスカラーであることは命題\ref{prop:finite-square-ratio}の証明による。
また$gU=U$なら$gQ=Q$なので$g$もスカラーである。従って符号の数は表示通りである。
前節で求めた$U$での接次元は0であり、各符号倍はスキームの自己同型なので全点が被約である。
最後に単元倍は第二窓の積の核を変えず、$U$自身は対称完全交叉からAS延長を持つ。
定理\ref{thm:finite-square-chart}と系\ref{cor:two-window-extension}で一意性を得る。
\end{proof}

\begin{theorem}[一般の第一窓に対する非延長]\label{thm:generic-finite-window-obstruction}
四次元の$E,F$を固定する。$\Gr(6,E\otimes F)$の稠密開集合上で、
第一関係空間$R$は四つの点対応が全て有限な標準AS対象に延びない。
より詳しくは、前定理の第一窓を含む開集合で、有限枝の第二窓は基底変更を除き
$R^{\mathrm{op}}\subset F\otimes E$に限られる。ここで$\mathrm{op}$は二文字の交換を表す。
\end{theorem}
\begin{proof}
順序付き二十点の配置を動かし、一次式と二次式の評価階数がそれぞれ4と10の開集合を取る。
各点の斉次座標を局所的に自明化する。平方が$Q$に入る四空間の相対Grassmann閉スキームは基底上固有であり、
各符号倍はその切断となる。前定理のファイバーでは全点の相対接空間が零なので、射はそこで不分岐である。
固有性により不分岐でない部分の像を除ける。不分岐射の切断は開埋込みであり、ここでは分離性から閉でもある。
全切断の補の像も閉で、前定理の点を含まない。従ってその近傍では全解が符号切断で尽くされる。
異なる切断の一致する閉集合も除けばこの有限被覆の次数は$2^{19}$のままである。

関係空間のGrassmann多様体上で、六本の双一次式の共通零点が被約な長さ20の完全交叉となる条件は開である。
実際、例の近傍で普遍零点は射影的で相対次元零、各零点で六式の相対Jacobi行列は可逆なので有限エタールである。
点を順序付ける有限エタール被覆と各行の自明化の上で、第一射影へ先の開条件を引き戻す。
結論は行の順序と単元倍に不変なので第一窓の開近傍を得る。
有限枝の延長があれば定理\ref{thm:finite-square-chart}から第三空間は第一空間の単元倍である。
従って第二の積の核は$R^{\mathrm{op}}$となる。

$A=R\otimes E\subset E\otimes F\otimes E$とし、両端の交換を$\sigma$と書く。
必要な三次交叉は$A\cap\sigma A$であり、標準AS対象では四次元である。
ところが一般の$R$でこの交叉は零となる。検算ファイルの整数係数の六列では、
$A$と$\sigma A$の挿入行列を並べた$64\times48$行列が標数101で階数48を持つので、
同じ小行列式は標数零でも非零である。階数48の開集合と先の非空開近傍は
Grassmann多様体の既約性で稠密開集合に交わる。これが最初の主張を与える。
\end{proof}

\begin{proposition}[反転する二窓の障害と二周期性]\label{prop:finite-flip-obstructions}
第一窓の点対応が有限とし、第二窓が$R^{\mathrm{op}}$であるとする。
三次交叉は両端を交換する対合の正固有空間に入り、その次元は
\[
 R\otimes E\longrightarrow \bigwedge^2 E\otimes F,\qquad
 (e\otimes f)\otimes e'\longmapsto(e\wedge e')\otimes f
\]
（単純テンソルからの制限）の余階数である。この行列は$24\times24$であり、
AS延長があれば階数は20、第三階数は$\rho\le10$となる。

さらに前定理の符号倍による復元が両射影に成立する被約有限枝では、任意のAS延長の周期は2を割る。
この場合、以下の$36\times21$行列の核はちょうど一次元で、
その生成元は六次元空間$R$上の係数行列として可逆である。
\[
 \operatorname{Sym}^2 R\ \longrightarrow\
 \bigwedge^2 E\otimes\bigwedge^2 F .
\]
写像は$R\subset E\otimes F$を用いた二組の交代化で、$r\otimes r$では
$4\times4$行列$r$の二次小行列式を並べる規約を取る。
\end{proposition}
\begin{proof}
両端交代な三次交叉の元は、$R$の基底六式の、第一文字の一次式を係数とする多項式恒等式を与える。
この一次関係は零である。実際、六式の有限完全交叉のKoszul分解を$(2,1)$でひねると、
外次数2から5の項の直線束$\mathcal O(2-j,1-j)$は全コホモロジーが零であり、
外次数6の非零項は$H^6$だけである。従って外次数1の大域切断から
$H^0(\mathcal O(2,1))$への写像は単射である。
交叉は対合で保たれ、正固有部分はちょうど表示した交代化の核である。
AS延長の標準テンソルはこの三次交叉と第四文字のテンソル積に入るので、
反対位置$(0,2)$では$\operatorname{Sym}^2 E$に入り、階数は10以下となる。

両射影の復元を順に使うと最初の四文字の座標空間を
$E,F,E,F$と同定でき、最初の三つの関係は$R,R^{\rm op},R$となる。
ここで第二の適用では最初の反転窓の第一射影が元の第二射影である。
先の交代部分の消滅を続く窓にも使えば標準テンソル$w$は
二つの端点交換$(0,2),(1,3)$のそれぞれで不変である。
従って第四の関係も$R^{\rm op}$、標準四周期表示の関係列は二周期となる。
さらに$w\in R\otimes R$は二組を交換しても不変なので$\operatorname{Sym}^2 R$に入り、
隣接切断階数6よりその係数行列は可逆である。
標数零の分解
\[
 \operatorname{Sym}^2(E\otimes F)=
 (\operatorname{Sym}^2 E\otimes\operatorname{Sym}^2 F)
 \oplus(\bigwedge^2 E\otimes\bigwedge^2 F)
\]
より、表示した写像の核は両方の端点交換で不変なテンソルからなる。

核の任意の元$w'$は第一と第二の関係条件を満たすので、
最後の文字を縮約して得る写像は同じ四次元三次交叉に入る。
$w$の対応する写像は同型だから、最後の$F$にある$g\in\operatorname{End}(F)$を
作用させて$w'$を得る。$w'$の最後の二文字は$R$に入り、$w$の隣接切断は全階数なので
$(1\otimes g)R\subset R$である。
各零点対$(x,y)$について$(x,g^*y)$も全関係で消える。$x$上の核は直線のため
$y\wedge g^*y=0$である。第二射影では二次式の評価が単射なのでこれは多項式恒等式となり、
$g$はスカラーである。従って核は$kw$で尽くされる。
二十点例では二行列の階数はそれぞれ20と20で、後者の核の係数行列は可逆である。
一般の階数条件からAS正則性を結論するにはなお全次数条件が要る。
\end{proof}
